\section{Experiments}
\label{sec:experiments}

We evaluate \toolname{} as a \emph{selective} detector: it may act only on pages
whose evidence is trustworthy enough, and should abstain otherwise. The selective
behavior is the contribution, so the seven questions below probe \emph{when} the
detector should be trusted, how robust that behavior is to its hyperparameters, and
\emph{what} it returns when it acts, rather than raw labelling accuracy:
\begin{compactitem}
\item[(\textbf{RQ1})] Selective detection (\autoref{sec:exp_rq1}): does grounded
evidence-agreement give a better risk--coverage trade-off than label-level
reliability, without sacrificing full-coverage detection accuracy?
\item[(\textbf{RQ2})] Right-label, wrong-reason (\autoref{sec:exp_rq2}): does
\toolname{} flag the pages a panel labels correctly but for the wrong reason --
the pages the reliability baselines act on confidently?
\item[(\textbf{RQ3})] Component attribution (\autoref{sec:exp_rq3}): how much do
evidence-level agreement, groundedness, and panel diversity each contribute?
\item[(\textbf{RQ4})] Calibration (\autoref{sec:exp_rq4}): is the trust score
calibrated, so its reported probability of correctness can be acted on directly?
\item[(\textbf{RQ5})] Sensitivity (\autoref{sec:exp_rq5}): how does the selective
trade-off respond to the two hyperparameters -- panel size and consensus threshold --
and is the default robust?
\item[(\textbf{RQ6})] Explanation (\autoref{sec:exp_rq6}): what does a returned
explanation look like, and does it make the act/abstain decision auditable to an
operator?
\item[(\textbf{RQ7})] Adversarial robustness (\autoref{sec:exp_adv}): when an
attacker corrupts a cited cue -- cloaking the form action or removing the brand logo --
does grounded evidence-agreement \emph{fail safe} (groundedness drops and the detector
abstains) where label-confidence baselines stay overconfident?
\end{compactitem}

\subsection{Setup}
\label{sec:exp_setup}

\sstitle{Benchmark and labels}
We evaluate on \textsc{PhishSel}, a $1{,}330$-page benchmark ($998$ held-out test) built
from the Phishpedia corpus~\cite{lin2021phishpedia}, whose phishing pages were originally
collected from public feeds~\cite{openphish,phishtank}. To make the task
\emph{legitimacy} rather than \emph{page type}, we pair the phishing pages with legitimate
\emph{credential-collection} (login) pages rather than benign homepages -- the latter are
several times larger than the minimal phishing pages and let a model separate the classes
by size alone (see \autoref{sec:limitations}). The phishing half is brand-balanced and
split evenly into easy (unrelated-domain) and hard (brand-in-domain lookalike) pages, the
latter forming the RQ2 subset. Each page carries its URL, DOM, and rendered screenshot;
the corpus does not include TLS certificates or redirect chains, so those two cue types are
inactive and the active schema is \{brand, form-action-domain, credential-intent,
logo-brand\} plus a derived brand--domain consistency check. Labels are available at
collection. We use a class-balanced calibration split ($332$) and a disjoint, balanced test
split ($998$), fit all thresholds and the calibrator on the former, and report every number
on the latter (\autoref{tab:dataset}).
\begin{table}[!h]
\caption{\textsc{PhishSel} composition. Phishing pages are drawn from the Phishpedia
corpus~\cite{lin2021phishpedia} (feeds collected mid-2019--late-2020), brand-balanced and
split evenly into easy (unrelated-domain) and hard (brand-in-domain look-alike) pages
($266$ easy / $233$ hard on the test split); benign pages are credential-collection (login) pages, so
the comparison turns on legitimacy rather than page type. The calibration split fits all
thresholds and the isotonic calibrator; the held-out test split carries every reported
number.}
\label{tab:dataset}
\footnotesize
\setlength{\tabcolsep}{6pt}
\begin{tabular}{lccc}
\toprule
Split & \#pages & \#phishing & \#benign \\
\midrule
Calibration      & \phantom{0,}332 & \phantom{0,}166 & \phantom{0,}166 \\
Test (held-out)  & \phantom{0,}998 & \phantom{0,}499 & \phantom{0,}499 \\
\midrule
Total            & 1{,}330 & \phantom{0,}665 & \phantom{0,}665 \\
\bottomrule
\end{tabular}
\end{table}


\sstitle{Agents, tools, and baselines}
The default panel is $M{=}3$ deliberately diverse \emph{evidence agents}: GPT-4o reading
the screenshot as the vision-language agent, and two different-family text LLMs run
locally -- Llama-3.2-3B (Meta) and Qwen2.5-3B (Alibaba) -- each able to call the
verification tools of \autoref{sec:ground}. The two text agents are deliberately from
different families: RQ3 shows this cross-family diversity, not the agent count, is what
carries the result (\autoref{tab:panel_models} lists the exact models and decoding
settings). All methods receive the same captured page snapshot and no
label at inference time, so any difference reflects how each method scores trust
rather than what information it sees. \toolname{} emits constrained evidence cues
in the schema of \autoref{sec:elicit}, normalizes values as in \autoref{sec:agree},
and verifies only the consensus cues before calibration. We compare against six
reliability baselines, each of which shares this panel or a member of it so that the
comparison isolates \emph{where} agreement is measured:
\begin{compactitem}
\item \emph{(B1) Single-model confidence}~\cite{guo2017calibration}: the strongest
single model with temperature-scaled softmax confidence.
\item \emph{(B2) Verbalized confidence}: the same model asked to state its
confidence.
\item \emph{(B3) Self-consistency}~\cite{wang2023selfconsistency}: the same strongest single
model (GPT-4o), sampled $k{=}5$ times at temperature $0.7$, with the majority-verdict
agreement as the score; run on the full $998$-page test split ($k\times$ API cost).
\item \emph{(B4) Majority vote}~\cite{ensembleornot2024}: panel label agreement.
\item \emph{(B5) Label-agreement}~\cite{dawid1979maximum}: reliability-weighted
label agreement over the panel.
\item \emph{(B6) MultiPhishGuard}~\cite{multiphishguard2025}: a separate GPT-4o
consolidator reads the panel's per-agent verdicts and cited evidence and outputs a
deliberated verdict and confidence, used as the score.
\end{compactitem}
To position \toolname{} against the dedicated phishing literature we also report the
LLM brand-and-domain checker \emph{(D3) PhishLLM}~\cite{tan2024phishllm}, re-run as a
standalone detector on \textsc{PhishSel}. The reference-based matcher \emph{(D1)
Phishpedia}~\cite{lin2021phishpedia} is re-run in a containerized \texttt{detectron2} on the
full test split (\autoref{tab:detect}); \emph{(D2) PhishIntention}~\cite{liu2022phishintention} layers a
heavier credential-page and interaction stage on the same detector core, so we cite its
reported numbers rather than re-running it. In our main runs the brand cue
is grounded by a page-content brand check; we additionally show (\autoref{sec:exp_rq1}) that
a reference detector (D1) can be substituted as this verifier with no loss of ranking, so
the panel \emph{subsumes} its signal directly. Methods that require
\emph{live} counterfactual interaction with the target -- notably DynaPhish's online
reference expansion~\cite{dynaphish2023} -- are out of scope, because that interaction
cannot be replayed on archived snapshots.
\begin{table}[!h]
\caption{The $M{=}3$ evidence-agent panel. The two text agents are deliberately from
different model families and run locally via Ollama on a $16$\,GB laptop (no GPU); the
vision agent is GPT-4o through its API. All agents decode deterministically (temperature
$0$), so the only stochastic choice in the pipeline is the calibration/test partition.}
\label{tab:panel_models}
\footnotesize
\setlength{\tabcolsep}{3.0pt}
\begin{tabular*}{\columnwidth}{@{\extracolsep{\fill}}llllc@{}}
\toprule
Agent & Model & Family & Modality & Host \\
\midrule
$A_1$ & Llama-3.2-3B & Meta    & text   & Ollama (local) \\
$A_2$ & Qwen2.5-3B   & Alibaba & text   & Ollama (local) \\
$A_3$ & GPT-4o       & OpenAI  & vision & API \\
\bottomrule
\end{tabular*}

\footnotesize{Text agents: $8{,}192$-token context, temperature $0$. Vision agent:
\texttt{detail=low}, temperature $0$. The MultiPhishGuard consolidator (B6) and the
verbalized/single-model baselines (B1, B2) reuse GPT-4o.}
\end{table}


\sstitle{Tasks and metrics}
Every method produces a per-page trust score, and we evaluate that score as a
selective detector via \eqref{eq:objective}. The selective metrics are the area
under the risk--coverage curve (\emph{AURC}, selective risk integrated over
coverage); selective accuracy and false-positive rate at $80\%$ coverage
(\emph{SelAcc$_{80}$}, \emph{FPR$_{80}$}); the coverage attainable at $99\%$
selective accuracy (\emph{Cov$_{99}$}), the operating point an operator cares about
most; and the expected calibration error of the trust score (\emph{ECE}). For RQ2 we
add the \emph{flip rate}: the fraction of pages whose panel label changes under a
small evidence-targeted perturbation (a logo morph~\cite{hao2024logomorph} or a
form-action cloak), which measures reason fragility. Because the specialized
detectors carry no abstention mechanism, we also report standard full-coverage
detection quality -- accuracy, precision, recall, and F1 with every page acted on --
so they can be compared at all. We define each abbreviation here and use it
unchanged in every table and figure below.

\sstitle{Protocol}
Operating points are set without touching the test pages: the thresholds for
SelAcc$_{80}$ and Cov$_{99}$ and the isotonic calibrator are all fit on the
calibration split and then applied unchanged to the held-out test split. For RQ2,
each perturbation targets the cited evidence cue while preserving the page's true
label, so a flip measures reason fragility, not a changed ground truth.

\sstitle{Statistical reporting}
Numeric cells report mean$_{\pm\text{95\% CI}}$ from a paired bootstrap over pages
(all methods share the same page draws). At $998$ test pages the intervals are wide, so
we describe a difference as significant only when the intervals do not overlap and as
suggestive otherwise; we are explicit about which is which in the text. The proposed
method is \textbf{bold} and the best baseline is \underline{underlined} in every table.
As the panel runs at temperature $0$ (deterministic given a page), the only stochastic
choice is the calibration/test partition; to confirm the headline is not an artifact of one
draw, we re-partition the $1{,}330$ pages five ways (stratified $75/25$) and refit the
calibrator each time. The metrics are stable across splits -- AURC $1.8\pm0.2$,
SelAcc$_{80}$ $97.5\pm0.3$, ECE $0.018\pm0.006$, Cov$_{99}$ $0.29\pm0.04$ -- and on AURC,
selective accuracy, and coverage the split reported in \autoref{tab:main} sits among the more
conservative draws (its ECE, $0.015$, is marginally better than the cross-split mean).

\sstitle{Run environment and cost}
The two text agents run locally on a 16\,GB commodity laptop via Ollama (no GPU) at
temperature $0$; the vision agent and the consolidator are GPT-4o through its API; the
DOM, redirect, logo (CLIP), and calibration code are CPU-only. Every model call is cached,
so the comparison is reproducible and re-scoring variants (the ablations) needs no
re-running. We log prompts, raw agent outputs, value-normalized cues, tool results,
thresholds, and per-page trust scores, and release the schemas, prompts, tool wrappers,
calibration code, and figure-generation scripts. \toolname{} has no training phase: per
page it runs the $M{=}3$ agents plus the cheap tool checks. \autoref{tab:cost} gives the
per-page budget: the panel's API cost is dominated by a single GPT-4o vision call
($\approx\$8$ per $1{,}000$ pages), since the two text agents run locally for free, so
\toolname{} fits a commodity $16$\,GB laptop with no GPU. The $M{=}3$ panel is a fixed
$3\times$ model-call overhead over a single-model detector (\autoref{fig:pareto}), amortized
because only the abstaining fraction $1-\phi(\thr)$ is escalated. At that budget \toolname{} stays
level with the panel baselines while reaching markedly higher deployable coverage,
and it dominates the five-sample self-consistency baseline (B3, \autoref{tab:main}: AURC
$5.1$ vs.\ $2.3$, and Cov$_{99}{=}0$ -- it never reaches the $99\%$-accuracy operating point)
at one-fifth its API cost, since B3 issues five GPT-4o vision calls per page where \toolname{}
issues one. Resampling one strong model improves selective accuracy over single-call B1/B2 but
does not cure its overconfidence: on the visually convincing brand replicas that dominate the
corpus, the model is confidently wrong and the five samples agree on the wrong \textsf{benign}
verdict, which is precisely the label-level trust that grounded agreement is designed to
replace.
\begin{table}[!h]
\caption{Per-page compute budget. The panel uses two local 3B text agents (free, on the
$16$\,GB laptop) and one GPT-4o vision call, so API cost is dominated by that single call;
the two text agents add no API cost. Dollar cost is at GPT-4o list price with detail-low
images; latency is indicative wall-clock on the CPU-only laptop (providers' API latency is
not controlled by us). The $M{=}3$ panel is a fixed $3\times$ model-call overhead over a
single-model detector, amortized by abstaining on the $1-\phi(\thr)$ fraction.}
\label{tab:cost}
\footnotesize
\setlength{\tabcolsep}{5pt}
\begin{tabular}{lccc}
\toprule
Method & Calls/page & API \$/1k & Latency/page \\
\midrule
\textbf{\toolname{}} (panel)    & 3 (2 local\,+\,1 API) & $\approx$8 & $\approx$20\,s \\
\;\,+\,B6 consolidator          & +1 API                & +$\approx$3 & +$\approx$3\,s \\
D3 PhishLLM                     & 1 API                 & $\approx$5 & $\approx$4\,s \\
D1 Phishpedia                   & 0 (local CPU)         & 0          & $\approx$15\,s \\
B1/B2 single-model              & 1                     & $\approx$2 & $\approx$3\,s \\
B3 Self-consistency ($k{=}5$)   & 5 API (same model)    & $\approx$40 & $\approx$15\,s \\
\bottomrule
\end{tabular}
\end{table}

% Generated PNG. Source: figures/gen_experiment_figs.py
% Regenerate: uv run --with matplotlib --with numpy python figures/gen_experiment_figs.py
\begin{figure}[!h]
\centering
\includegraphics[width=\linewidth]{fig_pareto}
\caption{Cost--quality trade-off: per-page model calls (analytical) versus AURC
($\downarrow$). At the $M{=}3$ panel's call budget \toolname{} has the lowest AURC in point
estimate (on par with B6 within CI),
below the other three-call panel methods (B4/B5/B6) and far below the single-call
confidence read-outs.}
\Description{A scatter plot of AURC against per-page model calls. PhishProof sits at
the bottom of the cost-quality frontier: at three model calls it has lower AURC than
the label-agreement, majority-vote, and MultiPhishGuard baselines at the same cost,
and far lower than the single-call confidence baselines.}
\label{fig:pareto}
\end{figure}


\sstitle{Reproducibility}
The pipeline is deterministic given the page snapshots and the dataset partition. We pin the
exact artifacts so the headline numbers can be regenerated: the panel is
\autoref{tab:panel_models} (Llama-3.2-3B and Qwen2.5-3B via Ollama, GPT-4o via API, all at
temperature $0$); the calibration and test splits are released with SHA-256 checksums
(\texttt{calibration.jsonl} \texttt{9f2c9bb8\dots}, \texttt{test.jsonl} \texttt{391fc594\dots});
the code runs on CPython $3.14$ with CPU-only DOM, CLIP, and calibration code. Every model
call is content-hash cached, so all ablations and the design-justification probes
(\autoref{sec:exp_negatives}) re-score from cache without re-running any agent. We release
the cue schema and prompts, the tool wrappers, the calibrator and threshold-fitting code,
and the figure- and table-generation scripts. Code, manifests, and headline result files are
available at \url{https://github.com/phamhyta/phishproof}; page snapshots and the full agent
cache are archived separately (Zenodo DOI upon acceptance).

\sstitle{Tool soundness}
\autoref{prop:hall} assumes the tools are sound (\autoref{ass:sound}); we check this for
\emph{every} active verifier, not only the brand check, by scoring each tool as a detector of
an independently-derivable target property (\autoref{tab:verifier_soundness}). The two
deterministic structural verifiers are reliable: the form-action verifier reproduces the
DOM-parsed action domains exactly (recall and specificity $1.00$), and the credential
verifier never misses a real password field (recall $1.00$). The two reference/perceptual
verifiers are accurate but soft: the brand-grounding tool recovers the impersonated gold
brand on $87.6\%$ of pages and spuriously grounds an absent brand on only $5\%$ ($95\%$
specificity), while the CLIP logo tool matches the gold brand at $89.0\%$ but grounds a
\emph{wrong} brand $34\%$ of the time (specificity $0.66$). Two soundness gaps are thus
explicit: the credential verifier over-fires on hint-only fields (precision $0.62$), and the
CLIP logo verifier is a soft matcher. Both are consistent with \autoref{prop:hall}
\emph{bounding} rather than eliminating gauge error, and the residual recall/precision gaps
are one reason groundedness is conservative on this corpus.
\begin{table}[!h]
\caption{Per-cue-type verifier soundness. Each grounding tool is scored as a detector of an
independently-derivable target property on the 998-page test split: the brand and logo
verifiers against the dataset brand annotation (specificity uses a random wrong brand), the
form-action verifier against the DOM-parsed action domains, and the credential verifier
against the strict presence of an \texttt{<input type=password>} field. The deterministic
structural verifiers are exact (form-action) or high-recall (credential), while the
perceptual CLIP logo verifier is a \emph{soft} matcher (specificity $0.66$). The
credential verifier never misses (recall $1.0$) but over-fires on hint-only fields
(precision $0.62$).}
\label{tab:verifier_soundness}
\footnotesize
\setlength{\tabcolsep}{3.0pt}
\begin{tabular*}{\columnwidth}{@{\extracolsep{\fill}}llcccc@{}}
\toprule
Cue type & Verifier & Prec. & Recall & Spec. & $n$ \\
\midrule
\texttt{brand\_claim}        & HTML-brand   & ---   & 0.876 & 0.950 & 499 \\
\texttt{form\_action\_domain} & DOM form     & ---   & 1.000 & 1.000 & 1067 \\
\texttt{credential\_intent}   & DOM cred.    & 0.624 & 1.000 & ---   & 998 \\
\texttt{logo\_brand}          & CLIP ViT-B/32 & ---  & 0.890 & 0.660 & 200 \\
\bottomrule
\end{tabular*}

\footnotesize{Specificity $=1-$ false-grounding rate against a random wrong value. The
logo verifier is subsampled to 200 pages with a detected logo box; others use the full
applicable split.}
\end{table}


\subsection{Selective Detection}
\label{sec:exp_rq1}

\sstitle{Full-coverage detection accuracy}
To answer \textbf{RQ1} we first establish the floor: at full coverage \toolname{} is a
solid detector, so its trust layer sits on top of -- not in place of -- accurate
labelling. \autoref{tab:detect} reports full-coverage quality with every page acted on.
\toolname{} reaches $91.9\%$ accuracy and $91.3$ F1 at very high precision ($98.8\%$),
ahead of the re-run LLM checker D3 ($83.7\%$ / $81.6$); but \toolname{} is an $M{=}3$
panel under a majority vote while D3 is a single model, so this margin reflects
ensembling as much as the trust layer, and we make no detector-leaderboard claim. The
re-run reference matcher D1 (Phishpedia, $n{=}998$) trails on F1 ($76.6$)
at high precision ($94.2$): it fires only on confident logo matches against its
reference list and abstains on the rest, so its recall ($64.5$) is bounded by brand
coverage; D2 adds a heavier stage on the same core and is cited. Beyond a competitor, D1
also serves as a \emph{grounding tool}: substituting it for the page-content brand check
leaves the selective ranking essentially unchanged (AURC $2.33$ vs.\ $2.29$ on the full
$998$ pages), so the panel can invoke a specialized detector to ground its brand cue
rather than re-derive it -- the ``subsumes the detector'' property of \autoref{sec:method},
measured here for D1 only (D3 substitution as a grounding tool is left to future work).
The intended reading is
only that detection is not sacrificed for the abstention and explanation the rest of the
section measures.
\begin{table}[!h]
\caption{Full-coverage detection quality on \textsc{PhishSel} (no abstention).
\toolname{} (the $M{=}3$ panel under a majority vote) and D1--D3 are re-run on the
full $998$-page test split (mean$_{\pm\text{95\% CI}}$, paired bootstrap).
D1 Phishpedia uses a containerized \texttt{detectron2} reference matcher ($\sim$15\,s/page on CPU).
D2 PhishIntention adds a heavier credential-page and interaction stage on the same detector
core, so we cite its reported results. \toolname{} leads the re-run reference detector D1
and the LLM checker D3 on accuracy and F1, but it is a three-model panel against a single
LLM checker, so the margin reflects ensembling as much as the trust layer; the point is that
detection is not sacrificed for the abstention and explanation the rest of the section
measures.}
\label{tab:detect}
\footnotesize
\setlength{\tabcolsep}{4.0pt}
\begin{tabular}{lcccc}
\toprule
Detector & Acc$\uparrow$ & Prec.$\uparrow$ & Rec.$\uparrow$ & F1$\uparrow$ \\
\midrule
D1 Phishpedia~\cite{lin2021phishpedia}        & 80.3$_{\pm2.6}$ & 94.2$_{\pm2.4}$ & 64.5$_{\pm4.3}$ & 76.6$_{\pm3.3}$ \\
D2 PhishIntention~\cite{liu2022phishintention} & \multicolumn{4}{c}{cited (heavier CRP/interaction stage); see \cite{liu2022phishintention}} \\
D3 PhishLLM~\cite{tan2024phishllm}             & 83.7$_{\pm2.4}$ & 93.3$_{\pm2.6}$ & 72.5$_{\pm4.1}$ & 81.6$_{\pm2.8}$ \\
\midrule
\textbf{\toolname{}} (full cov., $n{=}998$)   & \textbf{91.9}$_{\pm1.7}$ & \textbf{98.8}$_{\pm1.0}$ & 84.8$_{\pm2.5}$ & \textbf{91.3}$_{\pm1.7}$ \\
\bottomrule
\end{tabular}
\end{table}


\sstitle{Risk--coverage and deployment coverage}
Scoring agreement on grounded evidence, rather than on labels, gives the best selective
trade-off on the headline metrics. \autoref{tab:main} reports the comparison and
\autoref{fig:riskcoverage} the curves: \toolname{} has the lowest AURC ($2.3$, a
$38\%$ reduction over the label-agreement baseline B5), the highest selective accuracy
at $80\%$ coverage ($97.0\%$), and by a clear margin the best-calibrated trust score
(ECE $0.015$ versus $\ge0.026$ for every baseline). With $998$ test pages the bootstrap
intervals are wide, so the AURC lead over the strongest baseline (B6 MultiPhishGuard,
$3.0$) is suggestive rather than statistically separated (paired $\Delta$AURC $-0.4$, $95\%$
CI $[-1.5,+0.8]$, $p=0.44$). The direction, however, is robust to how the test set is
partitioned: \toolname's AURC is the lower of the two on $21/21$ $70\%$ subsamples and on
$4/5$ stratified folds, so the lead is consistent in sign even where it is underpowered to
separate. The calibration and selective-accuracy gaps are
unambiguous (paired $\Delta$ECE $-0.066$, CI $[-0.088,-0.043]$; $\Delta$SelAcc$_{80}$
$+2.4$, CI $[+1.0,+3.9]$). The deployment number Cov$_{99}$ is the one metric a baseline wins:
B6 reaches $0.36$ against \toolname's $0.25$, because grounded evidence is sparse on the
many correctly-labelled-but-weakly-grounded benign pages, which \toolname{} ranks low
and therefore abstains on. All panel-based methods meet at the shared full-coverage
panel error -- \toolname{} does not change the majority label, only the order in which
pages are trusted. The decisive comparison is against the label-aggregation baselines
(B4, B5): they carry a high ECE ($0.25$) because a label-vote share is not a calibrated
trust signal, whereas \toolname{} reads agreement on \emph{grounded evidence}, which is
where the calibration and AURC gains come from -- not merely from polling several models.
\begin{table}[!h]
\caption{RQ1. Selective detection on \textsc{PhishSel} (998 test pages). Lower
AURC/FPR/ECE and higher SelAcc/Cov are better. Cells are mean$_{\pm\text{95\% CI}}$
from a paired bootstrap over pages; \textbf{best}, \underline{best baseline}. The
$\Delta$ row is relative to B5. \toolname's gains in ECE and SelAcc$_{80}$ over the best
baseline are statistically significant (paired bootstrap CI excludes zero, marked
$\ast$); its AURC lead over B6 is the best point estimate but not statistically separated.
Self-consistency (B3) is the same strongest single model (GPT-4o) resampled $k{=}5$ times at
temperature $0.7$; its row is a point estimate ($k\times$ cost, no bootstrap CI).}
\label{tab:main}
\footnotesize
\setlength{\tabcolsep}{2.6pt}
\resizebox{\columnwidth}{!}{%
\begin{tabular}{lccccc}
\toprule
Method & AURC$\downarrow$ & SelAcc$_{80}\uparrow$ & FPR$_{80}\downarrow$ & Cov$_{99}\uparrow$ & ECE$\downarrow$ \\
\midrule
B1 Single-model conf.\         & 8.6$_{\pm2.4}$ & 92.4$_{\pm1.9}$ & 1.0$_{\pm0.9}$ & 0.00$_{\pm0.02}$ & 0.081$_{\pm0.017}$ \\
B2 Verbalized conf.\           & 5.7$_{\pm1.5}$ & 92.4$_{\pm1.8}$ & 0.9$_{\pm1.1}$ & 0.26$_{\pm0.03}$ & 0.026$_{\pm0.016}$ \\
B3 Self-consistency           & 5.1 & 96.0 & 0.4 & 0.00 & 0.041 \\
B4 Majority vote               & 8.5$_{\pm2.4}$ & 92.4$_{\pm1.9}$ & 1.0$_{\pm0.9}$ & 0.00$_{\pm0.02}$ & 0.252$_{\pm0.017}$ \\
B5 Label-agreement             & 3.7$_{\pm1.2}$ & 93.6$_{\pm1.8}$ & 1.5$_{\pm1.3}$ & 0.10$_{\pm0.23}$ & 0.238$_{\pm0.017}$ \\
B6 MultiPhishGuard             & \underline{3.0}$_{\pm1.1}$ & \underline{94.6}$_{\pm1.6}$ & 1.5$_{\pm1.3}$ & \underline{0.36}$_{\pm0.24}$ & \underline{0.083}$_{\pm0.016}$ \\
\textbf{\toolname{} (GEA)}     & \textbf{2.3}$_{\pm0.9}$ & \textbf{97.0}$^{\ast}_{\pm1.2}$ & \textbf{1.3}$_{\pm1.1}$ & \textbf{0.25}$_{\pm0.15}$ & \textbf{0.015}$^{\ast}_{\pm0.013}$ \\
\midrule
$\Delta$ vs B5                 & $-38\%$ & $+3.4$ & $-0.2$ & $+0.15$ & $-0.22$ \\
\bottomrule
\end{tabular}%
}
\end{table}

% Generated PNG. Source: figures/gen_experiment_figs.py
% Regenerate: uv run --with matplotlib --with numpy python figures/gen_experiment_figs.py
\begin{figure}[!h]
\centering
\includegraphics[width=\linewidth]{fig_riskcoverage}
\caption{Risk--coverage on \textsc{PhishSel} (RQ1, 998 pages). Lower selective risk is
better; the ring marks $80\%$ coverage. All panel methods meet at the shared
full-coverage error ($8.1\%$); \toolname{} changes only the abstention order and has
low risk as coverage tightens (best point-estimate AURC, on par with B6 within CI;
decisively the best calibration and selective accuracy).}
\Description{A line chart of selective risk versus coverage. PhishProof has the lowest
risk curve as coverage tightens and among the lowest area under the curve among the plotted
methods, which all meet at the shared full-coverage error.}
\label{fig:riskcoverage}
\end{figure}


\sstitle{Failure modes}
At full coverage \toolname{} errs on $81$ of $998$ pages, almost all in one direction:
$76$ are missed phishing pages and only $5$ are false alarms, matching the high-precision,
lower-recall profile of \autoref{tab:detect}. The misses are evidence-poor by construction:
$96\%$ carry no agreed brand cue and $66\%$ reach no consensus at all -- minimal or
hard-to-read pages on which the panel never converges. Crucially, the trust score
\emph{knows} this: errors carry a mean $\geascore$ of $0.20$ against $0.62$ on correct pages,
so the selective layer abstains on most of them and only $15$ of the $81$ errors are
\emph{confident} ($\geascore>0.5$). The dominant failure is thus a safe one -- abstention on
a page whose evidence does not cohere -- rather than a confident wrong verdict, which is the
behaviour the selective design is meant to produce. \autoref{tab:failure_modes} breaks the
errors down: the only genuinely dangerous mode -- a confident missed-phish that carries an
agreed brand cue -- accounts for just $3$ of $998$ pages.
\begin{table}[!h]
\caption{Failure-mode taxonomy of the $81$ full-coverage errors on the $998$-page test
split. Almost all errors are evidence-poor missed phishing pages on which the panel never
converges, and they carry a low mean $\geascore$ -- so the selective layer abstains on them.
Only $15$ of the $81$ errors are \emph{confident} ($\geascore{>}0.5$), and just $3$ are
confident missed-phish with an agreed brand cue (the only truly dangerous mode). Correct
pages average $\geascore{=}0.62$ for comparison.}
\label{tab:failure_modes}
\footnotesize
\setlength{\tabcolsep}{3.0pt}
\begin{tabular*}{\columnwidth}{@{\extracolsep{\fill}}lccc@{}}
\toprule
Failure mode & $n$ & mean $\geascore$ & confident ($>0.5$) \\
\midrule
FN: no consensus cue        & 50 & 0.00 & 0 \\
FN: cues but no brand       & 23 & 0.46 & 8 \\
FN: brand in consensus      & 3  & 0.81 & 3 \\
FP: false alarm             & 5  & 0.67 & 4 \\
\midrule
All errors                  & 81 & ---  & 15 \\
\bottomrule
\end{tabular*}
\end{table}


\subsection{Right-Label, Wrong-Reason Pages}
\label{sec:exp_rq2}
\textbf{RQ2} asks whether low GEA flags pages the panel labels correctly but for a
fragile reason -- the right-label, wrong-reason (RLWR) case. We test this with a small
evidence-targeted perturbation (a form-action cloak that rewrites the cited form to an
on-brand domain) that preserves the page's true label, and measure how often the panel
verdict flips. The effect is present but \emph{modest and benchmark-dependent}, and we
report it as such. On the unrestricted test split, where every phishing page sits on an
obviously off-brand host, verdicts are stable regardless of GEA and the gradient is
essentially flat. The RLWR effect only appears on the harder \emph{lookalike} subset
(brand-in-domain hosts such as \textsf{steamcoommuniity.com}; $233$ hard pages in test,
$n{=}200$ correctly labelled by the panel), and even there it is small:
\autoref{tab:rlwr} bins these correctly-labelled pages into GEA quintiles, and
fragility concentrates in the lowest quintile ($12\%$ flip) against $\le5\%$ above it.
Split coarsely by GEA level instead, the low-GEA group ($\geascore{<}0.5$, $n{=}98$) flips
$7.1\%$ against $6.0\%$ for the high-GEA group ($\geascore{>}0.72$, $n{=}67$) -- an overall
gap of only $1.2\times$, modest precisely because the effect is concentrated in the lowest
bin rather than spread across the range. So GEA
does rank the most fragile lookalike verdicts lowest -- which is the direction RQ2
predicts and the behaviour an operator wants from the abstention -- but the sharp
partition the illustrative design assumed does not materialise on this benchmark, whose
phishing pages carry few shared hallucinations for grounding to catch. We treat this as
a limitation of the benchmark rather than evidence against the mechanism
(\autoref{sec:limitations}). A natural \emph{evasive} subset isolates the effect more
clearly: filtering the $423$ correctly-labelled phish for pages with a minimal DOM
($<\!5$\,KB) or no detectable logo region yields $72$ pages, the closest the corpus
comes to logo-free/minimal-cue phishing. On this evasive subset the form-action cloak
flips $12.5\%$ of verdicts against $5.4\%$ on the rest ($2.3\times$), logo occlusion
$15.9\%$ vs.\ $8.3\%$ ($1.9\times$), and the white-box all-cue attack $19.4\%$ vs.\
$7.7\%$ ($2.5\times$) -- in every case the GEA-fragile pages are the ones a small
evidence-corruption flips. The mechanism is real; on a corpus richer in such pages the
flat-on-clean gradient would resolve into the sharp partition the design predicts.
\begin{table}[!h]
\caption{RQ2. Correctly labeled \emph{lookalike} phishing pages (brand-in-domain hard
subset: $n{=}200$ correctly labeled pages of $233$ hard pages in the test split) binned into
GEA quintiles. ``Flip'' is the label-flip rate under a
small evidence-targeted perturbation (form-action cloak) that preserves the true label.
Fragility is concentrated in the lowest-GEA bin ($12\%$ vs $\le5\%$ above it); a coarser
split by GEA \emph{level} -- the low group ($\geascore{<}0.5$, $n{=}98$) vs the high group
($\geascore{>}0.72$, $n{=}67$) -- gives a modest $7.1\%$ vs $6.0\%$ ($1.2\times$), small
because the effect is concentrated in the lowest bin rather than spread. On the
unrestricted test split, where the off-brand signal is unambiguous, the gradient is
flat. The strong right-label-wrong-reason effect is thus benchmark-dependent (see
Limitations), not the dramatic split the illustrative design assumed.}
\label{tab:rlwr}
\footnotesize
\setlength{\tabcolsep}{4.0pt}
\begin{tabular*}{\columnwidth}{@{\extracolsep{\fill}}lcc@{}}
\toprule
GEA quintile (low $\to$ high) & mean GEA & Flip$\downarrow$ \\
\midrule
Q1 (lowest)  & 0.29 & 12\% \\
Q2           & 0.50 & 2.5\% \\
Q3           & 0.55 & \phantom{0}5\% \\
Q4           & 0.72 & \phantom{0}5\% \\
Q5 (highest) & 0.92 & \phantom{0}5\% \\
\bottomrule
\end{tabular*}
\end{table}


\subsection{Component Ablation}
\label{sec:exp_rq3}
\emph{Panel diversity} is the load-bearing component: \autoref{tab:ablation} removes one
piece at a time. By far the largest effect is removing cross-family diversity -- swapping
the two distinct-family text agents for two same-family ones while holding the count at
$M{=}3$: on the matched subset this raises AURC from $1.3$ to $27.5$ and drops even
full-coverage accuracy from $0.90$ to $0.68$, because two same-family agents make
correlated errors that outvote the lone vision agent. This is why \emph{diversity}, not
panel size, is the design variable. Replacing evidence-level agreement with the label-vote
share costs the next most ($2.3\to3.5$ AURC, approaching the label-agreement baseline B5
($3.7$) of \autoref{tab:main}), confirming that the selective gain comes from scoring agreement on
evidence rather than on labels. Removing the calibrator leaves the \emph{ranking} intact
(AURC unchanged) but inflates the trust score's ECE from $0.015$ to $0.331$: calibration
governs the reported probability, not the order -- exactly the split of duties RQ4
examines. Dropping groundedness, by contrast, barely moves AURC ($2.3\to2.1$): on this
benchmark the agreed cues almost always survive their tool check, so there are few shared
hallucinations for $\grd$ to discard and grounding contributes little to the
\emph{ranking}. Its role is the guarantee of \autoref{prop:hall} and the verifiability of
the returned explanation rather than a measurable AURC gain here -- a point we return to
in \autoref{sec:limitations} and demonstrate directly under adversarial cue corruption in
\autoref{sec:exp_adv}, where groundedness is what drives the detector to abstain.
\begin{table}[!h]
\caption{RQ3. Component ablation. AURC/FPR$_{80}$ ($\times100$) and ECE; lower is
better. The first four rows are on the full 998-page test split; the calibrated ECE
of each row is reported (the $-$calibration row uses the raw, un-calibrated score).
The $-$diversity row is on a 100-page subset on which the full diverse panel scores
AURC $1.3$, so swapping the cross-family text agents for two same-family ones raises
AURC to $27.5$ and full-coverage accuracy falls $0.90\!\to\!0.68$ -- by far the
largest effect, which is why \emph{diversity}, not panel size, is the design variable.
The GPT-4o vision agent is held fixed in both arms, so the row isolates the effect of
cross-family \emph{text} diversity rather than confounding it with a change of vision model.}
\label{tab:ablation}
\footnotesize
\setlength{\tabcolsep}{3.0pt}
\begin{tabular*}{\columnwidth}{@{\extracolsep{\fill}}lccc@{}}
\toprule
Variant & AURC$\downarrow$ & FPR$_{80}\downarrow$ & ECE$\downarrow$ \\
\midrule
\textbf{\toolname{} (full)}                       & \textbf{2.3} & \textbf{1.3} & \textbf{0.015} \\
\;$-$ calibration                                  & 2.3 & 1.3 & \underline{0.331} \\
\;$-$ groundedness $\grd$ (\req{3})                & 2.1 & 1.2 & 0.014 \\
\;$-$ evidence-level $\agr$ (\req{2})              & \underline{3.5} & 1.3 & 0.015 \\
\;$-$ diversity (same-family)\textsuperscript{$\dagger$} & \underline{27.5} & \underline{48.6} & --- \\
\bottomrule
\end{tabular*}

\footnotesize{\textsuperscript{$\dagger$}100-page subset; the diverse panel scores
AURC $1.3$ on the same pages.}
\end{table}


\subsection{Why a Parameter-Free, Small-Model Design? Negative Results for Plausible Extensions}
\label{sec:exp_negatives}
The ablation shows which components are load-bearing; a natural follow-up is whether several
``obvious'' extensions improve on the deliberately simple design. We probe each -- four by
re-scoring the cached agent outputs at no extra cost, and one by re-running the panel with
larger models -- and find that the parameter-free product $\agr\!\cdot\!\grd$ over a fixed,
diverse, small-model panel is hard to beat.

\sstitle{A learned aggregator does not beat the crisp product}
The fusion $\geascore{=}\agr\!\cdot\!\grd$ has no learned parameters. We ask whether learning
a combiner over the per-cue-type signals $(a_t,g_t)$ -- which expose more structure than the
two scalars $\agr,\grd$ -- lowers AURC. \autoref{tab:hfgea} fits four combiners of increasing
capacity on the same $332$-page calibration split: a weighted geometric mean, a $2$-additive
Choquet integral, a per-type logistic model, and a Takagi--Sugeno--Kang fuzzy system. None
improves the headline ranking metric. The high-capacity fuzzy and Choquet models
\emph{overfit} the small calibration split (AURC $3.9$--$4.4$, Cov$_{99}$ collapsing to
$0.01$); the only competitive learned model, logistic regression, merely trades a small
AURC loss ($2.5\to2.9$) for a marginal Cov$_{99}$/FPR$_{80}$ gain. The crisp product remains
best on AURC ($2.5$) and ECE ($0.015$), so we keep it and read this as evidence that the
calibration-scale data does not support a learned fusion.
\begin{table}[!h]
\caption{Does a \emph{learned} aggregator beat the parameter-free $\agr\!\cdot\!\grd$?
Each combiner is fit on the 332-page calibration split over the same per-cue-type
$(a_t,g_t)$ signals (four active types) and evaluated on the 998-page test split.
AURC/SelAcc$_{80}$/FPR$_{80}$ are $\times100$. No learned combiner improves the headline
ranking metric (AURC): the high-capacity fuzzy/Choquet models overfit the small
calibration split (AURC $3.9$--$4.4$, Cov$_{99}$ collapses), and the only competitive
learned model (logistic) merely trades AURC for a small Cov$_{99}$/FPR gain. The crisp
product remains best on AURC and ECE, justifying the parameter-free design.}
\label{tab:hfgea}
\footnotesize
\setlength{\tabcolsep}{3.0pt}
\begin{tabular*}{\columnwidth}{@{\extracolsep{\fill}}lccccc@{}}
\toprule
Aggregator & AURC$\downarrow$ & SelAcc$_{80}\uparrow$ & FPR$_{80}\downarrow$ & Cov$_{99}\uparrow$ & ECE$\downarrow$ \\
\midrule
\textbf{crisp $\agr\!\cdot\!\grd$ + iso} & \textbf{2.47} & 96.87 & 1.27 & 0.210 & \textbf{0.015} \\
weighted geom. $\agr^a\grd^b$            & 2.56 & 97.12 & 1.20 & 0.210 & --- \\
Choquet 2-additive                       & 4.44 & 94.49 & 1.34 & 0.012 & 0.249 \\
logistic (per-type)                      & 2.85 & 97.12 & \textbf{0.97} & \textbf{0.228} & 0.020 \\
TSK fuzzy (heavy $L_2$)                  & 3.89 & 96.24 & 1.30 & 0.012 & 0.115 \\
\bottomrule
\end{tabular*}

\footnotesize{All aggregators are compared on identical per-cue-type features fit on the
calibration split, so the contrast within this table is controlled; the crisp baseline here
($2.47$) is the same $\agr\!\cdot\!\grd$ product re-scored over the per-type representation and
differs marginally from the deployed pipeline's headline AURC ($2.3$, \autoref{tab:main}). The
weighted-geometric mean is reported without the isotonic layer, hence its ECE is omitted.
Higher-capacity learned aggregators were regularized heavily and cross-validated; the
overfitting is intrinsic to the 332-page calibration scale.}
\end{table}


\sstitle{Per-class recalibration trades coverage for ranking}
The one metric a baseline wins is Cov$_{99}$ (\autoref{sec:exp_rq1}). We test the natural
fix -- calibrating predicted-phish and predicted-benign pages with separate isotonic curves.
It improves AURC ($2.5\to2.2$) and SelAcc$_{80}$ but \emph{lowers} Cov$_{99}$ ($0.21\to0.11$),
because the per-class curves introduce more tied probabilities at the high end, and Cov$_{99}$
rewards a dense cluster of distinct top-confidence pages. The Cov$_{99}$ gap to B6 is therefore
not a calibration-asymmetry artifact -- after calibration $99\%$ of correctly-detected
phishing pages already exceed trust $0.8$ -- but a property of how isotonic regression
discretizes the high-confidence tail. We keep the single calibrator, which is the better
operating-coverage choice.

\sstitle{The panel cannot be made cheaper by a cascade}
Because the vision call is the only paid one, a tempting deployment optimization is a cascade
that decides easy pages with the two free local text agents and escalates to GPT-4o only when
they are uncertain. This fails: the two small text agents agree on the verdict for only
$1$ of $998$ pages, and a text-only trust score is no better than random (AURC $\approx50$).
The cross-modal vision agent is load-bearing, which corroborates the diversity ablation
(\autoref{sec:exp_rq3}) -- the panel's power comes from agreement on \emph{grounded evidence}
across diverse modalities, not from label consensus among similar models, so the panel cannot
be thinned without losing the signal.

\sstitle{Memory-augmented brand recovery has headroom but low precision}
The dominant failure mode is missed phishing pages on which the panel never names the
impersonated brand ($73$ of $76$, \autoref{sec:exp_rq1}). We test an episodic prior that
retrieves the most DOM-similar calibration pages and borrows their brand, gated by grounding.
The headroom is real -- $32$ of the $76$ misses have a groundable retrieved brand -- but
converting the prior into a verdict is precision-sensitive: only at a strict similarity gate
does it become net-positive, recovering $4$ missed phishing pages with \emph{zero} new false
positives, while looser gates create more false alarms than they fix. High-precision brand
recovery is thus a promising but open direction rather than a free win (\autoref{sec:limitations}).

\sstitle{A stronger text panel does not help -- cue elicitation, not model scale, is the bottleneck}
Because the panel is $3$B-class to fit a $16$\,GB laptop, a natural objection is that larger
models would change the result. We swap both text agents for their $2$--$3\times$ larger
siblings (Llama-3.1-8B and Qwen2.5-7B, same prompts, same GPT-4o vision) and re-run on a
$120$-page stratified subset. The stronger panel does \emph{not} improve detection: AURC rises
$4.2\to9.5$ and accuracy falls $90.0\to76.7$. The cause is not malformed output (JSON validity
is $100\%$ for both) but \emph{cue elicitation}: the larger models return a verdict with an
\emph{empty} cue set far more often ($30$--$51\%$ of pages vs.\ $\approx0\%$ for the $3$B
Qwen), and a verdict with no typed cues grounds nothing, so $\geascore$ loses its signal. This
is consistent with the rest of this section: the method's power comes from agents emitting
\emph{groundable evidence}, which the compact instruction-tuned models do reliably, not from
raw model scale. The $3$B panel is therefore a principled choice, not only a hardware
concession.

\subsection{Trust-Score Calibration}
\label{sec:exp_rq4}
The calibrated trust score is faithful: \toolname{} attains the lowest expected
calibration error of any method, ECE $0.015$ (\autoref{tab:main}), against $0.081$ for
single-model confidence (B1) and $0.24$ for the label-aggregation baselines (B4, B5),
whose label-vote share is not a calibrated probability at all. So when \toolname{}
reports a trust near $0.9$, close to $90\%$ of those verdicts are in fact correct. This
faithfulness is what makes the operating point usable: because the calibrator is fit on
the held-out labeled split rather than guessed, the reported probability can be turned
directly into an act/abstain threshold at a target error budget, which is the premise the
selective results of RQ1 rely on. It is also the clearest separation between \toolname{}
and every baseline (\autoref{tab:main}), and the component the ablation confirms is
essential to it (\autoref{tab:ablation}, $-$calibration).

\subsection{Sensitivity Analysis}
\label{sec:exp_rq5}
Two hyperparameters govern the selective trade-off: the panel size $M$ and the consensus
threshold $k$. \autoref{tab:sensitivity} sweeps both by re-scoring the cached agent outputs,
so no agent is re-run. Panel size matters sharply. A single agent has no cross-agent
agreement signal -- $\agr$ is trivially $1$, so $\geascore$ collapses to bare groundedness --
and a two-agent panel turns the majority rule into a brittle unanimity; only at $M{=}3$,
where ``two of three'' is an achievable majority, does grounded evidence agreement become
well-defined (AURC $2.3$ versus $\ge35$ for $M{\le}2$), consistent with the diversity
ablation of \autoref{sec:exp_rq3}. The consensus threshold is likewise non-trivial: requiring
any single agent ($k{=}1$) admits uncorroborated cues, while a unanimous rule ($k{=}3$)
discards genuine evidence, so the strict-majority rule ($k{=}2$) is the sweet spot
(AURC $2.3$ versus $10.4$ and $6.0$). The reported configuration is therefore not incidental:
$M{=}3$ with majority consensus is the minimum at which the method is well-defined and at
which AURC is lowest.
\begin{table}[!h]
\caption{RQ5. Sensitivity to the two hyperparameters, by re-scoring the cached agent
outputs on the $998$-page test split (AURC $\times100\,\downarrow$, SelAcc$_{80}\,\uparrow$).
\emph{Top:} panel size $M$ ($M{<}3$ averaged over all single agents / pairs). \emph{Bottom:}
the consensus threshold $k$ at $M{=}3$. The reported configuration ($M{=}3$, majority
$k{=}2$) is the minimum at which grounded evidence agreement is well-defined and gives the
lowest AURC.}
\label{tab:sensitivity}
\footnotesize
\setlength{\tabcolsep}{6pt}
\begin{tabular*}{\columnwidth}{@{\extracolsep{\fill}}lcc@{}}
\toprule
Configuration & AURC$\downarrow$ & SelAcc$_{80}\uparrow$ \\
\midrule
\multicolumn{3}{@{}l}{\emph{Panel size $M$ (strict-majority consensus)}}\\
\quad $M{=}1$ (single agent)        & 35.2 & 65.2 \\
\quad $M{=}2$ (pair)                & 38.6 & 64.2 \\
\quad \textbf{$M{=}3$ (full panel)} & \textbf{2.3} & \textbf{97.0} \\
\addlinespace[2pt]
\multicolumn{3}{@{}l}{\emph{Consensus threshold $k$ (at $M{=}3$)}}\\
\quad $k{=}1$ (any agent)           & 10.4 & 91.4 \\
\quad \textbf{$k{=}2$ (majority)}   & \textbf{2.3} & \textbf{97.0} \\
\quad $k{=}3$ (unanimous)           & \phantom{0}6.0 & 93.2 \\
\bottomrule
\end{tabular*}
\end{table}


\subsection{Qualitative Study}
\label{sec:exp_rq6}
\textbf{RQ6} asks what \toolname{} hands an operator when it acts, and whether that
artifact makes the act/abstain decision auditable; \autoref{tab:qual} answers by walking
two real pages through the pipeline. On the impersonation page~(A), a page on
\texttt{turkmcse.com} posing as the IRS, the vision agent and one text agent converge on the
same brand claim, both read the off-brand credential form from the DOM, and every agreed cue
survives its tool check -- the CLIP logo check confirms the IRS mark and a rule flags that
the claimed brand does not match the host. \toolname{} acts and returns the agreed, verified
cues as the explanation, every line naming a tool the operator can re-run. Page~(B), a real
NatWest phishing page on the look-alike host \texttt{natwesttreats.co.uk}, shows the
abstention path: the panel \emph{splits} (one agent flags the form, one reads NatWest
branding, one sees nothing), so no cue reaches a majority and the single perceptual cue that
is relaxed in -- a claimed NatWest logo -- \emph{fails} its CLIP check ($0.03$). Groundedness
collapses ($\grd{=}0.03$) and with it $\geascore{=}0.01$, so \toolname{} abstains and
escalates rather than emit the panel's bare majority label -- which here would have been the
\emph{wrong} \textsf{benign} verdict, turning a silent false negative into a human review.
The two cases are the method in one table: the explanation is exactly the evidence the panel
agreed on and the page confirmed, and where no such evidence survives its check the detector
declines to explain and defers.
\begin{table*}[!t]
\caption{Qualitative study: two \emph{real} \textsc{PhishSel} pages and what \toolname{}
returns. A cue joins the explanation only if it reaches consensus ($\cons$: a strict
majority of agents cite it, or a single-source perceptual cue is relaxed in, raising
$\agr$) \emph{and} its tool re-derives it from the page (raising $\grd$); \cmark/\xmark{}
mark consensus membership.}
\label{tab:qual}
\footnotesize
\setlength{\tabcolsep}{4pt}
\begin{tabularx}{\textwidth}{@{}>{\raggedright\arraybackslash}X ccc@{}}
\toprule
Cited cue (type $=$ value) & Agents & Tool $\gver$ & $\cons$? \\
\midrule
\multicolumn{4}{@{}p{\textwidth}@{}}{\textbf{(A)} \texttt{turkmcse.com} (impersonates IRS) -- acted \textbf{phishing}, $\hat\rho{=}0.99$} \\
\quad brand $=$ IRS                         & 2/3   & detector \cmark & \cmark \\
\quad form-action $=$ \texttt{turkmcse.com} & 2/3   & DOM \cmark      & \cmark \\
\quad credential-intent $=$ yes             & 2/3   & DOM \cmark      & \cmark \\
\quad logo $=$ IRS                          & 1/3   & CLIP \cmark     & \cmark \\
\quad brand $\neq$ domain                   & der.\ & rule \cmark     & \cmark \\
\multicolumn{4}{@{}p{\textwidth}@{}}{\quad$\agr{=}1.00,\ \grd{=}1.00,\ \geascore{=}1.00$\,;\ \ explanation $=$ the five \cmark{} cues} \\
\addlinespace[2pt]
\multicolumn{4}{@{}p{\textwidth}@{}}{\textbf{(B)} \texttt{natwesttreats.co.uk} (real \textsf{phish}) -- \textbf{abstained} (panel split)} \\
\quad logo $=$ NatWest                      & 1/3 & CLIP $0.03$ \xmark & \cmark \\
\quad brand $=$ NatWest                     & 1/3 & detector \cmark    & \xmark \\
\quad form-action $=$ (empty)               & 1/3 & DOM \xmark         & \xmark \\
\multicolumn{4}{@{}p{\textwidth}@{}}{\quad$\agr{=}0.50$ ($1$ consensus of $2$ pooled cues; the empty form-action is not pooled)$,\ \grd{=}0.03 \Rightarrow \geascore{=}0.01$\,;\ \ escalate} \\
\bottomrule
\end{tabularx}
\end{table*}


\subsection{Adversarial Robustness}
\label{sec:exp_adv}
\textbf{RQ7} asks what happens when an attacker corrupts the very cue a verdict cites --
the regime the clean corpus does not contain (\autoref{sec:exp_rq3}) and the one in which
grounding is meant to matter. On $423$ correctly-detected phishing pages ($420$ for the
logo attack, where three carried no detectable logo region) we apply two
evidence-targeted perturbations that \emph{preserve} the true phishing label -- a
\emph{form-action cloak} that rewrites the credential form to an on-brand-looking domain,
and a \emph{logo occlusion} that overwrites the brand logo region, synthesizing the
logo-free/minimal-cue phishing the corpus lacks -- and re-run the panel.
\autoref{tab:adversarial} reports the result. Two findings stand out.

\emph{The panel is hard to evade, but \toolname{} fails safe when it is.} Single-cue
attacks flip only $6.6$--$9.7\%$ of verdicts: the diverse, multi-signal panel rarely
collapses because no one corrupted cue carries the decision. More importantly, on the
pages an attack does flip to the wrong \textsf{benign} verdict, \toolname{} abstains at
the deployed operating point on $39$ of $40$ under occlusion and on \emph{all} of them
under the form cloak ($28/28$) and the combined attack ($41/41$) -- $\ge97.5\%$ overall,
mean trust $\approx0.58$, below $\thr$ -- turning a silent false negative into an
escalation. The single-model confidence baseline B1, by contrast, reports mean confidence
$1.00$ on those same evaded pages -- wholly blind to the corruption, because it reads
trust at the label, which the attack leaves untouched. The MultiPhishGuard consolidator
B6, which inspects the agents' rationales, is more responsive (mean $\approx0.68$ on the
evaded pages) but still above \toolname's $\approx0.58$ there and, lacking a calibrated
abstention rule, does not decline them.

\emph{Grounding is the component that activates under attack.} The visual attack (logo
occlusion) drives groundedness down ($\grd$ $0.95\to0.75$) and abstention up ($69\to88\%$):
the occluded logo fails its CLIP re-derivation, so the cue the panel cites no longer
grounds and $\geascore$ falls. This is precisely the behaviour that is invisible on clean
data -- where $\grd$ barely moves the ranking (\autoref{sec:exp_rq3}) -- confirming that
tool-grounding earns its place as a fail-safe against corrupted evidence rather than as a
clean-corpus ranking gain. The form-action cloak alone moves $\grd$ little ($0.95\to0.90$),
because the cloaked form re-derives as on-brand and simply grounds a different cue; its
effect surfaces only on the evaded subset.

\emph{The fail-safe holds under a white-box, all-cue adversary.} The strongest test is an
adaptive attacker that knows every verifier and corrupts all of them at once -- cloaking the
form action, erasing the brand's textual identity from the DOM, and occluding the logo
(\autoref{tab:adversarial}, last row, $400$ pages). This is harder to ground than any single
attack ($\grd$ $0.95\to0.77$) and flips more verdicts ($12.5\%$), as expected when no cited cue
survives. Yet on \emph{every} one of the $50$ evaded pages \toolname{} abstains
($100\%$, mean trust $0.55<\thr$), while B1 again reports confidence $1.00$: the adaptive
adversary can suppress the evidence and force a wrong \textsf{benign} label, but cannot
manufacture \emph{grounded agreement}, so it cannot make \toolname{} act on that label. We do
not claim robustness to an adversary who controls the panel or the tools themselves
(\autoref{sec:related}); within cue-level corruption -- including this white-box, all-cue case
-- the detector declines to act rather than acting wrongly.
\begin{table}[t]
\caption{RQ7. Adversarial robustness on $423$ correctly-detected phishing pages ($420$ for
logo occlusion, three pages lacked a detectable logo) under
evidence-targeted perturbations that preserve the true (phishing) label. ``Evade'' is the
fraction whose panel verdict flips to \textsf{benign}; $\grd$ and ``Abstain'' are reported
\emph{clean}\,$\to$\,\emph{attacked} (Abstain is the share with $\geascore<\thr$ at the
deployed operating point). The last column is measured \emph{on the evaded pages only}:
how often \toolname{} abstains versus the single-model confidence baseline B1's mean
confidence. The visual attack (logo occlusion) drives groundedness down and abstention up,
and on evaded pages \toolname{} abstains on $97.5$--$100\%$ of them (logo occlusion
$97.5\%$; the other attacks $100\%$) while B1 stays maximally confident on the now-wrong
verdict. $^{\ddagger}$The adaptive (all-cue) attack is a white-box
adversary that corrupts \emph{every} active cue at once -- form-action cloak, brand-text
removal, and logo occlusion -- on $400$ pages; it flips more verdicts ($12.5\%$) but \toolname{}
still abstains on $100\%$ of the evaded pages.}
\label{tab:adversarial}
\footnotesize
\setlength{\tabcolsep}{3.6pt}
\resizebox{\columnwidth}{!}{%
\begin{tabular}{lccccc}
\toprule
 & & \multicolumn{2}{c}{clean\,$\to$\,attacked} & \multicolumn{2}{c}{on evaded pages} \\
\cmidrule(lr){3-4}\cmidrule(lr){5-6}
Attack & Evade$\downarrow$ & $\grd$ & Abstain & \toolname{} abst.\ & B1 conf.\ \\
\midrule
Form-action cloak & 6.6\% & .95$\to$.90 & 69$\to$60\% & 100\% & 1.00 \\
Logo occlusion    & 9.5\% & .95$\to$.75 & 69$\to$88\% & 97.5\% & 1.00 \\
Both              & 9.7\% & .95$\to$.80 & 69$\to$79\% & 100\% & 1.00 \\
Adaptive (all-cue)$^{\ddagger}$ & 12.5\% & .95$\to$.77 & 69$\to$83\% & 100\% & 1.00 \\
\bottomrule
\end{tabular}%
}
\end{table}


\subsection{Generalization to Unseen Brands}
\label{sec:exp_brands}
A single-corpus evaluation invites the question of whether the calibrated trust merely
memorizes the brands it was fit on. We test this with a \emph{leave-brands-out} protocol that
needs no new data: we split the brand vocabulary into two disjoint halves, fit the calibrator
and Cov$_{99}$ threshold on pages whose brand is in the \emph{seen} half, and evaluate on pages
whose brand is \emph{unseen} (the date-/brand-less benign pages are split at random). Averaged
over five brand partitions, each leaving $\approx700$ unseen-brand pages,
\autoref{tab:brands_out} shows the selective metrics on unseen brands match the
in-distribution split and on AURC slightly exceed it ($1.9\pm0.3$ vs.\ $2.3$), with calibration
holding (ECE $0.03$). Because the trust score reads agreement on \emph{grounded evidence} --
a brand impersonated on an off-brand domain grounds as inconsistent regardless of which brand
it is -- the operating point transfers to brands never seen during calibration. This is a
distribution-shift test \emph{across brands}, complementary to the time-ordered drift test
below; live cross-feed validation on freshly-collected pages remains future work
(\autoref{sec:limitations}).
\begin{table}[!h]
\caption{Leave-brands-out generalization. The brand vocabulary is split into two disjoint
halves; the calibrator and threshold are fit on pages whose brand is in the \emph{seen} half
and evaluated on pages whose brand is \emph{unseen} (benign pages split at random). Mean$\pm$std
over five brand partitions ($\approx$700 unseen-brand test pages each). The selective metrics
on entirely unseen brands match -- and on AURC slightly exceed -- the in-distribution split,
so the calibrated trust transfers across the brand distribution rather than memorizing brands.}
\label{tab:brands_out}
\footnotesize
\setlength{\tabcolsep}{3.0pt}
\begin{tabular*}{\columnwidth}{@{\extracolsep{\fill}}lccccc@{}}
\toprule
Split & AURC$\downarrow$ & SelAcc$_{80}\uparrow$ & FPR$_{80}\downarrow$ & Cov$_{99}\uparrow$ & ECE$\downarrow$ \\
\midrule
In-distribution (\autoref{tab:main}) & 2.3 & 97.0 & 1.3 & 0.25 & 0.015 \\
Unseen brands & $1.9_{\pm0.3}$ & $97.7_{\pm0.7}$ & $1.1_{\pm0.4}$ & $0.30_{\pm0.05}$ & $0.03_{\pm0.02}$ \\
\bottomrule
\end{tabular*}
\end{table}


\subsection{Robustness to Temporal Drift}
\label{sec:exp_drift}
A deployed detector is calibrated once and then meets pages from later, evolving campaigns,
so a trust score is only useful if it survives that gap. The \textsc{PhishSel} phishing pages
span roughly a year (2019-07 to 2020-10), which lets us test this directly: we fit the
calibrator and thresholds on a \emph{past} window ($\le$\,2020-07) and evaluate on the
disjoint \emph{future} window ($\ge$\,2020-08), with the date-less benign pages split at
random (\autoref{tab:drift}). Trained only on the past, \toolname{} holds its ranking and
selective accuracy on the future window (AURC $3.1$ vs.\ the in-distribution $2.8\pm0.6$;
SelAcc$_{80}$ $96.2$ vs.\ $96.9$), both within the in-distribution spread. Calibration is the
one quantity that drifts, and only slightly -- ECE rises from $0.021$ to $0.030$ -- yet it
stays an order of magnitude below the label-vote baselines' full-data ECE ($\ge0.24$,
\autoref{tab:main}). Because trust is read from \emph{grounded evidence} rather than a fitted
decision boundary, the operating point set on past data transfers to future pages without
re-tuning. This time-ordered protocol follows the temporal-bias discipline of
TESSERACT~\cite{pendlebury2019tesseract} -- calibration strictly precedes test in time --
and the stability we observe is consistent with \emph{agreement-on-the-line}: cross-model
agreement tracks accuracy under distribution shift~\cite{baek2022agreement}, which is why
a score read from agreement transfers where a fitted boundary need not. We note the gap is
modest: the corpus spans only mid-2019 to late-2020, so this probes resilience to short-horizon
drift, not to the multi-year evolution a deployed detector eventually meets
(\autoref{sec:limitations}).
\begin{table}[!h]
\caption{Robustness to temporal drift. The calibrator and operating thresholds are fit on a
\emph{past} window (phishing pages dated $\le$\,2020-07, $n{=}275$) and evaluated on the
disjoint \emph{future} window ($\ge$\,2020-08, $n{=}390$); benign pages, which carry no
campaign date, are split at random. The in-distribution row averages five temporally-mixed
$50/50$ splits ($\pm$ standard deviation). \toolname{} keeps its ranking and selective
accuracy across the $\sim$1-year gap; calibration degrades only slightly and stays far below
any baseline's full-data ECE ($\ge0.026$, \autoref{tab:main}).}
\label{tab:drift}
\footnotesize
\setlength{\tabcolsep}{5pt}
\begin{tabular}{lcccc}
\toprule
Split & AURC$\downarrow$ & SelAcc$_{80}\uparrow$ & Cov$_{99}\uparrow$ & ECE$\downarrow$ \\
\midrule
In-distribution (mixed)    & 2.8$_{\pm0.6}$ & 96.9$_{\pm0.6}$ & 0.26$_{\pm0.08}$ & 0.021$_{\pm0.007}$ \\
Past $\rightarrow$ future  & 3.1            & 96.2            & 0.23            & 0.030 \\
\bottomrule
\end{tabular}
\end{table}

